
\section*{Installment 1: 	extit{Folkebladet}, April 24, 1918}

\subsection*{The Congregation}

\textsc{I.}

It is not so difficult to be editor of \textit{Folkebladet} when friends are so ready to write as they have been of late. In any case, he need not write much himself. The difficulty is rather to find room for all the reading matter that is sent to the paper. From time to time, however, the editor may feel prompted to have a word in the matter, especially when it concerns the paper’s editorial position.

Now Dr. Evjen has, in a series of articles on ``the Position,'' thoroughly and clearly shown what, in his opinion, the Free Church’s doctrine of the congregation is, as set forth by Professors Sverdrup and Oftedal, and as it has found expression in \textit{Guiding Principles}. Although \textit{Folkebladet} has not written much editorally on that matter, there are nevertheless a couple of incidental statements that are weighed and tested, and found not to be sufficiently clear; yet no major objection is raised against them.

A more exact explanation or account from the side of \textit{Folkebladet} is perhaps not necessary or demanded; yet we shall try to write a little more directly on the subject that is here being discussed.

For the sake of clarity, we shall state our position in the following points:

\begin{enumerate}
\item In accordance with \textit{Guiding Principles}, we believe that the congregation is the proper form of the Kingdom of God on earth.
\item The church bodies are human constructions, which may indeed have significance for the advancement of the Kingdom of God, but are nevertheless not the proper form of the Kingdom of God.
\item We can agree to call \textit{Guiding Principles} essentially polity, that the Free Church’s polity is congregationalist, and that it differs from that of other Lutheran bodies.
\item The Free Church’s chief task, however, is not to labor for polity, but for the salvation of souls, for spirit and life.
\end{enumerate}

So far as we understand Dr. Evjen, this is also his position. His articles, however, have this time been chiefly referential; he has cited the statements of others and does not set forth so much his own personal opinion, or what he himself considers to be most in agreement with the Word of God. We understand that he has altered his view somewhat of late, and he hints, as his opinion, that if Prof. Sverdrup had lived long enough that he could have read Rudolf Sohm’s work on church law from 1914, this might perhaps have moved him to moderate his view with regard to principle no. 1. He further says that principle no. 1 requires our special attention, and he raises these questions: ``Is it too much to ask whether it is fully in accord with the Word? Is it absolutely excluded that the authors could have been mistaken? The opposition to this principle at any rate warrants the posing of such questions. Thus: questions that await an answer, tasks that await a solution.''

We assume, however, that Dr. Evjen still assents to this principle, and then we are agreed in the main matter. And the thought suggests itself involuntarily: Why not look more to the fundamental principles on which we are agreed, and less to the unessential things which are of such a kind that there may very well be room for divided opinions? We have such need to work together rather than to split apart.

Take, for example, the thorny question of the Free Church’s distinct position among the Norwegian Lutherans. Some hold that there is only a difference of polity; others that there is more, namely also a difference in spiritual direction. Some lay the weight upon polity, others upon spirit and life. And it might look as though the two parties stood so far from one another, indeed as though they were even the very opposite of one another. Perhaps, when it comes down to it, they are in reality nearer to one another than they themselves suppose.

It seems to us that the difficulty lies in the question of what is meant by polity. If we could obtain a sound definition of this concept, it would surely help us to understand one another better. Most people no doubt mean by polity the same as constitution, laws, rules, statutes for the government of a congregation, a church body, or a country---that is, the outward form of government. But then there are those who mean not only the forms, the rules, but also the spirit, the inward pulsating life, life with its forms (for life is presupposed).

It will readily be seen that there is a considerable difference between the two ways of viewing the matter, and herein lies the difficulty, we believe.

It may be that it is true that some of us have distinguished too much. But that too is not so simple. It appears that the elder Sverdrup, whose pupils we all believe ourselves to be, also distinguished between spirit and polity, for example in the following statement:

\begin{quote}
Without spirit there is no freedom in church or congregation. To suppose that one can produce freedom in the congregation by laws and constitutions is an error or delusion of the same kind as if a man should imagine that he could obtain steam power, if only he had a steam boiler, or that he could obtain wheat without end, if only he had a threshing machine. A good constitution and good laws are useful, if freedom is there; but if freedom is not in the congregation, then the law cannot produce it. Law does not create freedom; but if it is the right kind of law, then it serves to give freedom occasion to work and can even serve to preserve and protect freedom.

In the Norwegian state church there is at present quite great freedom for lay activity, because the law protects it. There is far, far greater freedom for lay activity than in the great majority of congregations here among us. That shows that freedom comes from spirit, not from church polity. It also shows that the law can serve freedom by protecting it; and conversely that good laws can become altogether ineffectual if there is no spirit in the congregation.

The first chief condition for a liberated congregation is therefore spirit, God’s Spirit. And if then the polity is sufficiently spacious, so that not all spiritual life is driven out of the congregation, then freedom too can accomplish its work.

If there is spirit in the congregation, then the first great effect of it will be a living zeal for the salvation of souls. That is the first and most important of all expressions of life and movements of freedom in the congregation, the very burning zeal of divine love for the salvation of souls.
\end{quote}

\noindent(\textit{Collected Writings} II, 21--22).

In this connection we also recall what Sverdrup once said in a lecture concerning the Congregationalists in this country. He called attention to the fact that they are the continuation of the old Puritans. They have preserved the congregationalist polity, but the spirit is not the same; for the Congregationalists have now become the most rationalistic of all the Reformed church bodies.

We may learn from this development among the Congregationalists. It would be fatal for us to settle ourselves down and think that we have grasped it, that we have reached the goal, because we have a free polity. Above all, it is a matter of having spirit, life. All that concerns laboring for awakening and for the preservation of spiritual life always remains the main matter. If there was anything our old leaders laid particular weight upon, it was precisely this.

Do not misunderstand us, as though we meant to insinuate that those who so strongly emphasize this matter of polity are opposed to spiritual life. Not at all. On the contrary, it is the purpose of this article to advance the claim that at bottom we want the same thing: that those who stress polity also want spirit and life, and that those who lay the weight on spirit and life also want a free polity---most of them at least. Why then devour one another?

That there are some who in some respect desire a different polity may well be. But if we in writing, speech, and conduct show Christian forbearance and the right brotherly spirit, those who would rob us of the foundation will not gain much following. We must make room for one another. There must be freedom in the Free Church. We are not served by having anyone set himself up as a dictator with a claim to sole right to decide what is truly Free Church.
That there are some who in some respect desire a different polity may well be. But if we in writing, speech, and conduct show Christian forbearance and the right brotherly spirit, those who would rob us of the foundation will not gain much following. We must make room for one another. There must be freedom in the Free Church. We are not served by having anyone set himself up as a dictator with a claim to sole right to decide what is truly Free Church.

We who speak so much about putting spiritual gifts to use in the practical work of the congregation ought also to be willing to let people, both pastors and laypeople, have opportunity to express themselves in our paper on questions of principle, even if they have not had occasion to read so many learned works. It is, after all, in the practical work of the congregation that the principles are to be carried through, and what is learned in the practical school of life is not to be despised. Besides, we all have access to the Bible---the Book---upon which our guiding principles are built; and when laypeople have the right and calling to expound the Word of God, they should surely also be given the right and opportunity to speak on the principles, for these are not holier than the Word itself.

As we have often stated, we are in favor of free, substantive discussion of questions of principle. We are not among those who would have everyone keep silent, and have no discussion at all. We believe in freedom of speech, and we believe that mutual, brotherly criticism does good. We all have our peculiarities, sharp edges, and faults that may need to be corrected. If criticism falls silent, everything will grow stagnant. We therefore do not want any strict censorship as in the days of absolutism. We want full freedom of speech, but indeed under responsibility.

We have spoken against personal attacks upon brothers, and we do so still; but we do not simply call it a personal attack when one answers another’s opinion, and perhaps even uses strong expressions when they are needed. No, let the ideas contend. Nor ought any of us to be so thin-skinned that we cannot bear criticism. To grow sulky, or to play the martyr, when one meets contradiction, is altogether too childish.

\textit{(To be continued.)}

